% Bản nháp chương 9 theo dàn bài mới, viết với bộ mẫu ở style/mau/.
% ===========================================================================
\chapter{Kết luận}
\label{sec:conclusion}

Chương~\ref{sec:problem} đặt câu hỏi gồm hai vế: huấn luyện mô hình định giá là
giải một bài toán tối ưu hóa, vậy giải bằng cách nào và giải tới mức nào thì
dừng. Chương này trả lời cả hai, và mỗi câu trả lời đi kèm điều kiện mà nó không
còn đúng, vì mọi kết luận trong báo cáo đều đo trên một bài toán bậc hai với
$d = 280$ và $\kappa = \num{268.31}$.

\section{Giải tới mức nào}
\label{sec:conclusion-how-far}

Sai số $10^{-6}$ là đủ, và từ chương~\ref{sec:enough} thì con số ấy có nghĩa cụ
thể thay vì là một quy ước: dưới mức đó, sai lệch giá so với nghiệm đóng không
quá một đơn vị tiền trên một chiếc máy giá trung vị. Ngưỡng có căn cứ đo được là
$1{,}35 \cdot 10^{-6}$, tức quy ước quen dùng chỉ chặt hơn 1,35 lần.

Kết luận đảo chiều theo dung sai chứ không theo thuật toán. Nếu chấp nhận sai
lệch mười đơn vị tiền, ngưỡng nới ra thành $4{,}0 \cdot 10^{-5}$, và khi đó SGD
bậc thang từ chỗ không đạt trở thành đạt sau 21 lượt duyệt dữ liệu. Dung sai của
bài toán ứng dụng, chứ không phải bản thân thuật toán, quyết định phương pháp nào
được coi là dùng được.

Chiều ngược lại cũng cần nói. Chạy chính xác hơn mức cần không mua thêm gì: từ
$10^{-8}$ trở xuống, sai lệch giá còn dưới \num{0.02} đơn vị tiền, và gradient
descent dừng ở $4{,}9 \cdot 10^{-8}$ còn cho RMSE thấp hơn nghiệm đóng một chút
vì $\wstar$ cực tiểu hàm mục tiêu trên tập huấn luyện chứ không trên tập kiểm
tra.

\section{Giải bằng cách nào}
\label{sec:conclusion-how}

Trên bài toán này, câu trả lời là giải thẳng. Bước Newton chính là nghiệm đóng
theo \eqref{eq:newton-one-step}, nên toàn bộ lời giải mất \SI{0.133}{\second},
nhanh hơn 3,5 lần so với thời gian gradient descent cần chỉ để chạm ngưỡng, và
bộ giải trực tiếp của scikit-learn làm đúng việc ấy trong \SI{0.235}{\second}.
Phương pháp lặp không có lý do tồn tại ở quy mô này.

Điều kiện đảo chiều là quy mô. Newton tốn $\bigO(nd^{2} + d^{3})$ mỗi vòng và
phải lập ma trận $d \times d$; khi $d$ lớn tới mức không lập nổi ma trận ấy thì
L-BFGS là phương pháp duy nhất trong báo cáo còn chạy được, vì nó chỉ giữ $m$ cặp
hiệu và tốn $\bigO(md)$ ngoài một lần lấy gradient. Ở $d = 280$ thì Newton-CG
chậm hơn Newton 10 lần, nhưng nó tồn tại đúng vì lý do ấy và cân bằng sẽ đảo ở
$d$ đủ lớn.

Điều kiện đảo chiều thứ hai là dữ liệu. Cả bảy cấu hình trong báo cáo đều giả
định toàn bộ ma trận thiết kế nằm trong bộ nhớ. Với dữ liệu đến theo dòng hoặc
lớn hơn bộ nhớ, SGD lấy lại chỗ đứng không phải vì nó hội tụ tốt hơn mà vì các
phương pháp còn lại không chạy được.

\section{Chọn tham số}
\label{sec:conclusion-parameters}

Ba lời khuyên rút từ chương~\ref{sec:mechanisms}, mỗi lời kèm bằng chứng và giới
hạn.

Chọn độ dài bước sát ngưỡng $2/L$ nhưng đừng chạm vào nó. Chênh 10\% quanh ngưỡng
đủ để lật từ hội tụ sau 20 vòng sang phân kỳ với $f$ lên $5{,}4 \cdot 10^{11}$,
và đúng tại ngưỡng thì một thành phần đứng yên nên lần chạy dừng ở
$4{,}9 \cdot 10^{-8}$. Lời khuyên này đòi biết $L$; khi không biết thì
backtracking thay thế được với cái giá là số lần đánh giá hàm, và cái giá ấy do
$\beta$ quyết định chứ không do $\alpha$, từ \num{1.40} lần thêm mỗi vòng với
$\beta = 0.5$ lên \num{4.1} lần với $\beta = 0.9$.

Giữ momentum nhất quán với độ dài bước thực tế. Công thức
$\beta = (\sqrt{\kappa}-1)/(\sqrt{\kappa}+1)$ chỉ đúng ở $t = 1/L$, và ghép nó
với bước do line search chọn làm hàm mục tiêu bùng lên $1{,}47 \cdot 10^{4}$.
Dùng \eqref{eq:momentum-from-step} theo bước thực tế, và thêm vào đó điều kiện
Armijo phải lấy $\alpha = 1/2$ để trùng bổ đề giảm \eqref{eq:descent-lemma}, hoặc
phải khởi tạo $t_0 = 1/L$.

Cho mỗi kích thước lô một độ dài bước riêng. Hằng số trơn của mất mát một mẫu lớn
hơn $L$ tới 1159 lần, nên một bước an toàn cho gradient đầy đủ làm SGD với lô một
mẫu phân kỳ. Và nếu muốn sai số tiếp tục giảm thì phải để $\eta_k$ tắt dần, vì
bước hằng chỉ đưa nghiệm tới một lân cận quanh $\wstar$; giảm theo bậc thang cho
kết quả tốt hơn bước hằng 165 lần với cùng ngân sách.

\section{Lý thuyết dùng được vào việc gì}
\label{sec:conclusion-theory}

Cận co rút \eqref{eq:gd-rate} mô tả đúng phần đuôi tới năm chữ số thập phân nhưng
lệch 93 lần ở giai đoạn đầu, vì phổ Hessian dồn về phía đầu lớn: 269 trong 280
trị riêng nằm trên $L/10$ và chỉ 8 trị nằm dưới $10\mu$. Hệ quả thực dụng là dùng
cận để ước lượng chạy thêm bao lâu thì được, dùng nó để lập ngân sách vòng lặp từ
đầu thì lãng phí gần hai bậc, và dùng nó để chọn cấu hình thì sai, vì momentum
hằng tính từ $\kappa$ vốn tối ưu theo lý thuyết lại chậm hơn momentum tăng dần
trên bài toán này.

Chỗ cận có ích nhất hóa ra là làm phép kiểm tra tính đúng đắn. Hệ số co rút quan
sát được nằm trên cận là dấu hiệu lỗi cài đặt chứ không phải bài toán khó, vì cận
đúng cho mọi hàm trong lớp, và nhóm đã dùng đúng phép kiểm ấy khi truy lỗi
momentum.

\section{Kết luận nào là sản phẩm của dạng bậc hai}
\label{sec:conclusion-quadratic}

Một phần kết quả trong báo cáo gắn với việc hàm mục tiêu là hàm bậc hai và sẽ đổi
trên bài toán phi tuyến, nên cần tách ra.

Chuyện Newton hội tụ sau đúng một vòng lặp là hệ quả trực tiếp của
\eqref{eq:derivatives}, rằng Hessian không phụ thuộc $w$, nên xấp xỉ bậc hai mà
Newton dựng không phải xấp xỉ. Trên hàm phi tuyến, Hessian đổi theo $w$, Newton
cần nhiều vòng và số vòng lặp trở lại thành đại lượng đáng so. Cùng lý do ấy làm
backtracking cho Newton không có việc gì để làm ở đây, vì bước đầy đủ luôn thỏa
điều kiện Armijo.

Việc phân rã lại Hessian mỗi vòng hay dùng lại phân rã cũ cho cùng một thời gian
cũng là sản phẩm của dạng bậc hai, vì cả hai chỉ chạy một vòng. Trên bài toán phi
tuyến đây là hai lựa chọn khác hẳn nhau về chi phí.

Ngược lại, ba cơ chế ở chương~\ref{sec:mechanisms} không phụ thuộc dạng bậc hai.
Ngưỡng ổn định theo hằng số trơn, phương sai của gradient ngẫu nhiên, và chi phí
mỗi vòng lặp đều là những đại lượng có mặt ở mọi bài toán trơn, chỉ khác là $L$
và $\mu$ khi ấy phải ước lượng địa phương thay vì đọc từ phổ của một ma trận cố
định.

\section{Tỉ lệ đáng nhớ nhất}
\label{sec:conclusion-scale}

Nếu chỉ giữ lại một con số từ báo cáo này thì nên giữ tỉ lệ sau. Mô hình định giá
sai trung bình \num{4245} đơn vị tiền trên một chiếc máy giá trung vị. Khoảng
cách giữa cấu hình tối ưu hóa tốt nhất và tệ nhất mà nhóm tự cài là \num{87.7}
đơn vị, tức khoảng 2\% của con số trên. Một quyết định về hệ số hiệu chỉnh thì
đắt hơn, \num{302} đơn vị, dù nó chỉ là một dòng trong phần chuẩn bị dữ liệu.

Nói cách khác, trên bài toán này toàn bộ phần tối ưu hóa quyết định khoảng 2\%
sai số cuối cùng, còn 98\% còn lại thuộc về việc chọn mô hình và chuẩn bị dữ
liệu. Điều đó không làm phần tối ưu hóa bớt đáng học, vì nó là phần quyết định
mất \SI{0.133}{\second} hay \SI{40}{\second} để có mô hình, và vì nó là phần duy
nhất trong toàn bộ quy trình có thể sai theo kiểu không nhận ra được: một cấu
hình phân kỳ thì thấy ngay, còn một cấu hình dừng ở $8{,}2 \cdot 10^{-4}$ vẫn trả
về một bộ trọng số trông hoàn toàn bình thường và chỉ đắt thêm \num{87.7} đơn vị
tiền mỗi máy.

\section{Phạm vi của các kết luận}
\label{sec:conclusion-scope}

Bốn giới hạn cần ghi rõ. Mọi con số thời gian đo trên một máy duy nhất, và
chênh lệch giữa các lần chạy của cùng một cấu hình lên tới 61\% ở trường hợp xấu
nhất, nên chỉ những khoảng cách lớn hơn mức đó mới nên đọc thành thứ hạng. Thứ
hạng giữa các thuật toán chỉ đo ở một giá trị $\kappa$, và
mục~\ref{sec:lambda-scope} nói rõ vì sao đó vẫn là dự đoán. Hàm mục tiêu là hàm
bậc hai, với các hệ quả đã tách ở mục~\ref{sec:conclusion-quadratic}. Và toàn bộ
báo cáo làm việc ở $d = 280$, tức ở quy mô mà bộ giải trực tiếp còn thắng, nên nó
không nói được gì về vùng mà phương pháp lặp thực sự cần thiết.
